\section{系统总体设计}

本章从整体层面阐述系统的设计理念、整体框架、功能组件划分以及实验组织方式。

\subsection{总体设计思路}

Thought2Text 数据集的规模为约 12000 条 trial，覆盖 40 个视觉类别，其数据量仅相当于主流视觉-语言模型预训练集的百万分之一左右。同时，该数据集缺少人工编写的自然语言图像描述，唯一的原始标注信息仅为类别标签。受限于上述条件，直接从 EEG 信号端到端地学习自由文本生成面临严重的瓶颈。针对这一问题，本设计采取由简单到复杂、从定量判别到生成式输出的渐进式策略，将完整系统划分为四个逐层递进的阶段。各阶段均具有清晰的输入输出约定和可独立验证的中间产物，前一阶段的成果可以直接被后续阶段复用或作为基础构件。

\textbf{第一阶段：CLIP 语义空间构建}。借助 CLIP 在海量图文数据上习得的跨模态对齐空间，搭建 EEG 与视觉语义之间的桥梁。具体而言，利用 CLIP 文本编码器将 40 个类别的文本 prompt 预先计算为固定的文本原型向量，同时通过 CLIP ViT 对整个图像集进行特征预缓存。该阶段为后续所有涉及 EEG 与图像语义交互的工作提供了统一的向量空间，避免了在有限 EEG 样本上重新学习语义结构的必要性。

\textbf{第二阶段：A2 EEG 语义融合}。本阶段设计并实现了 A2 temporal-spectral-spatial EEG encoder，用于从 EEG 信号中同时提取时域、频域和空域三方面的语义信息。在 CLIP 向量空间中，采用门控残差融合（gated residual fusion）机制将 EEG embedding 与退化图像的视觉 embedding 相结合，以分类交叉熵损失作为主导训练信号，并引入 shuffled EEG（打乱脑电）和 random EEG 两类对照实验来验证 EEG 语义贡献的真实性。

\textbf{第三阶段：token-level EEG 视觉增强}。第二阶段中的语义融合在 pooled embedding 层面进行，此时整幅图像的视觉信息已被压缩为单一向量。为了让 EEG 信息参与到生成式 VLM 的逐 token 处理流程中，本阶段引入 VTF 模块，在 CLIP ViT 输出的 visual patch token 层级注入 EEG token，利用 cross-attention 机制生成 EEG-enhanced vision tokens。

\textbf{第四阶段：生成式 VLM 集成与验证}。将上一阶段得到的 enhanced vision tokens 送入 Qwen2-VL-2B-Instruct 生成式模型，采用 LoRA 进行参数高效微调。通过对比不同生成方案与训练目标，并运用 best-of-N reranking 策略优化输出质量，最终构建起从退化图像和 EEG 信号到自然语言描述的完整生成通路。

\subsection{系统总体架构}

系统的整体架构如图~\ref{fig:overall-system} 所示。

\placeholderfigure{系统总体架构图}{overall-system-placeholder}{4.5cm}

数据在系统中的流转过程如下：退化图像首先经过冻结的 CLIP ViT 编码器，分别输出 pooled image embedding（512 维）和 visual patch tokens（维度为 $50\times512$）；与此同时，EEG trial 经 A2 encoder 处理，输出 pooled EEG embedding（512 维）和 EEG tokens（维度为 $4\times512$）。在语义预测分支中，两个 pooled embedding 经门控融合后与文本原型向量进行相似度计算，输出类别预测结果。在 token 级融合分支中，visual tokens 和 EEG tokens 通过 VTF 中的 cross-attention 产生 enhanced vision tokens。在生成式分支中，增强后的视觉 token 先经由本设计实现的轻量 prefix adapter 映射至 Qwen2-VL 可接受的嵌入空间，再与语义提示文本共同送入生成式解码器，经 LoRA 微调后生成自然语言描述。

\subsection{模块划分}

整个系统包含七个核心模块。表~\ref{tab:modules} 汇总了各模块的名称与其功能。

\begin{table}[htbp]
  \centering
  \caption{系统模块划分及功能说明}
  \label{tab:modules}
  \begin{tabularx}{\linewidth}{l X}
    \toprule
    模块名称 & 功能说明 \\
    \midrule
    图像编码模块 & CLIP ViT-B/32，提取 pooled embedding 和 visual patch tokens，模型参数冻结 \\
    EEG 编码模块 & A2 temporal-spectral-spatial encoder，提取 pooled EEG embedding 和 EEG tokens \\
    文本原型模块 & CLIP text encoder 预计算 40 类文本原型向量，训练期间冻结 \\
    语义融合模块 & Gated residual fusion 将图像与 EEG embedding 融合，与文本原型匹配进行分类 \\
    Token 融合模块 & VTF cross-attention 将 EEG tokens 融入 visual tokens，生成增强视觉 token \\
    生成式描述模块 & Qwen2-VL-2B-Instruct + LoRA，接收增强视觉 token 并生成自由文本描述 \\
    评估模块 & 多模式对照评测，统计各项指标，挖掘定性示例 \\
    \bottomrule
  \end{tabularx}
\end{table}

\subsection{实验流程}

图~\ref{fig:experiment-flow} 展示了从数据准备到最终评估的完整实施步骤。

\placeholderfigure{实验流程图}{experiment-flow-placeholder}{4.0cm}

整个流程包含以下环节：(1) 数据准备阶段——下载 EEG 数据、建立与 ImageNet 图像的软链接、构造 JSONL manifest 文件；(2) CLIP 特征预缓存——预先计算文本原型向量以及全部图像的视觉特征；(3) A2 EEG encoder 训练——以 InfoNCE 对比损失为目标进行 EEG 到 CLIP 空间的对齐训练；(4) 语义融合训练——联合训练门控融合与分类模块，在六种退化条件下执行多模式评估；(5) VTF token 融合训练——训练基于 cross-attention 的 token 层级融合模块；(6) 生成式模型训练——将 enhanced vision tokens 送入 Qwen2-VL，利用 LoRA 进行参数高效微调，对比多种训练目标与 LoRA 配置方案；(7) 全量测试与 reranking——在测试集上遍历所有组合进行评估，并基于最优 checkpoint 实施 best-of-N reranking；(8) 报告生成——汇总各项指标、筛选典型定性样例。

所有实验在单张 NVIDIA RTX 6000 Ada（48 GB 显存）上完成，采用 bfloat16 混合精度训练，未涉及多 GPU 分布式策略或 4-bit 量化技术。
